\subsection{One-bit OR proofs and their causal security model}\label{sec:mv:bit}
The three-move Schnorr OR proof of Chapter~10 has a challenge
in $\mathbb Z_P$. A generic simulator that guesses this entire
challenge succeeds with probability $1/P$, which is unsuitable
when $P$ is cryptographically large. We use the classical
\emph{one-bit challenge and rewinding} technique
\cite{Schnorr1991,CramerDamgardSchoenmakers1994}
to obtain an explicit malicious-verifier simulator.
A textbook treatment of the bit-challenge technique is
given in~\cite{Goldreich2001Foundations}.

\begin{definition}[One-bit OR mini-proof]\label{def:mv:mini}
Fix a public group commitment $C$, a nonempty finite set
$S=\{t_1,\dots,t_d\}\subset\mathbb Z_P$, and a witness
$(j_*,r)$ satisfying $C=g^{t_{j_*}}h^r$.
The prover executes Definition~\ref{def:realize:or},
except that the \emph{global} verifier challenge is
a bit $e\in\{0,1\}\subset\mathbb Z_P$.
The internal branch challenges $e_j$ and responses $z_j$
remain arbitrary elements of $\mathbb Z_P$ and must obey
$\sum_je_j=e\pmod P$.
An honest verifier samples the challenge bit uniformly
\emph{after} receiving the entire first-message vector $T$.
\end{definition}

\begin{lemma}[Perfect completeness of the mini-proof]\label{lem:mv:mini-complete}
For either global challenge bit, every honest execution
of Definition~\ref{def:mv:mini} satisfies the OR verifier equations.
\end{lemma}
\begin{proof}
The completeness proof of
Lemma~\ref{lem:realize:or-complete} is an identity
for every $e\in\mathbb Z_P$, hence in particular
for $e=0$ and $e=1$.
\end{proof}

\begin{definition}[Schedule-respecting malicious verifier]\label{def:mv:verifier}
A malicious verifier $V^*$ is a probabilistic
algorithm with a private random tape and a
rewindable state. At each OR mini-proof stage,
after observing the statement, previous transcript,
and first message $T$, it may either:
(i) send either bit $0$ or $1$;
or (ii) terminate with a distinguished abort symbol $\bot$.
Any other response is normalized to $\bot$.
The verifier may choose the bit as an arbitrary efficient
function of its entire visible history and its private tape.
All other public Ring Switching challenges may likewise
be chosen causally from the current visible history.
The verifier cannot demand an extraneous secret opening
or change the prescribed protocol equations.
\end{definition}

\begin{definition}[One-step black-box rewinding]\label{def:mv:rewind}
A simulator is allowed to save the verifier state
immediately before one OR first message, execute
the verifier on a trial first message, inspect the
bit challenge or abort, and restore the saved state
with the \emph{same} verifier random tape.
Only the current mini-proof is rewound.
After a trial is selected, the simulator continues
from the selected verifier state without changing
any earlier public message.
This is the standard classical black-box
rewinding model and does not grant rewinding
access to a quantum verifier.
\end{definition}

\begin{lemma}[A fixed statement admits a known witness in the honest case]\label{lem:mv:honest-witness}
For every bitness or parity OR statement generated
by an honest registration and a true affine relation
in Chapter~10, the honest prover knows a candidate
index $j_*$ and blinding $r$ required by
Definition~\ref{def:mv:mini}.
\end{lemma}
\begin{proof}
For bitness, the actual bit and its fresh Pedersen
blinding give $C=g^b h^r$ with candidate $b\in\{0,1\}$.
For an affine parity row, the true integer sum lies
in its permitted parity set by
Lemma~\ref{lem:realize:correct-candidate}.
Lemma~\ref{lem:realize:aggregate} supplies its known
aggregate blinding. Both are valid OR witnesses.
\end{proof}

\begin{remark}[Cost of choosing a shorter challenge]
The mini-proof has challenge entropy only one bit.
Its standalone knowledge error cannot be treated as
$1/P$. It must be repeated with fresh prover randomness
and fresh verifier challenge bits; the soundness analysis
of this repetition appears in
Section~\ref{sec:mv:forking}.
\end{remark}
